Multi-agent LLM debates suffer from majority bias: correct minority
hypotheses are eliminated prematurely before adequate evaluation, reducing
accuracy and collective reliability.
We propose \textbf{Adaptive Influence Balancing} (AIB), a framework that
monitors live process diagnostics and selectively intervenes to preserve
minority agents with high-quality reasoning.
We introduce the \emph{Premature Elimination Rate} (PER)—the fraction of
proposed correct hypotheses discarded before the final round—as a direct
measure of this failure mode.
Evaluating five conditions on 200 MMLU-Pro questions with Qwen2.5
(3~agents, 5~rounds), we find that diagnostic-triggered minority protection
achieves 60.0\% accuracy vs.\ 54.5\% baseline (+5.5~pp, $p\!<\!0.01$),
reduces PER by 61.3\% (8.2\%→3.2\%), and improves consensus reliability
to 60.6\%.
A learned predictor (Exp4) replaces hard thresholds with a calibrated
logistic model trained on process diagnostics and LLM-judged reasoning
quality; top predictors of minority correctness are explanation quality
($\beta\!=\!0.245$) and support trajectory ($\beta\!=\!0.179$).
Notably, adversarial loser resistance \emph{worsens} PER (+9.7\%) and
categorical minority protection paradoxically increases influence asymmetry,
demonstrating that trigger selectivity is the critical design choice.
Our results show that process-diagnostic control signals enable principled,
effective intervention in collective LLM reasoning.
